% =====================================================================
%  Chapter 22 — Capture
% =====================================================================

\chapter{Capture: The Solar System's Private Dark Matter}
\label{ch:capture}

The Solar System contains a private supply of dark matter, gravitationally
captured and concentrated, that exists nowhere else in the Solar System's
vicinity. It is the clearest sense in which the dark matter has been
``tapped'' in a way that no laboratory experiment has matched, and it is
also the sense that is least exploited.

\section{What capture does}Two things happen to a dark matter particle that
enters the Solar System's gravitational field.

It is \emph{accelerated}. As it falls towards the Sun it gains kinetic
energy, and it emerges from the encounter with a speed greater than the
local galactic value, at a rate that depends on the scattering radius. This
is the mechanism behind the dark matter wind and the hail.

It is \emph{captured}. A fraction of the incident particles are left on bound
orbits after the encounter, and those particles constitute a population that
stays in the Solar System. The captured population has a velocity
distribution concentrated at the escape velocity at the radius of the
scattering, which is low --- of order tens of kilometres per second near the
Earth --- and this is what makes it both a scientific target and a
phenomenological complication.

\begin{table}[H]
\centering
\caption{Approximate characteristic velocities for the dark matter populations
relevant to Solar System experiments. The bound population is more than an
order of magnitude slower than the halo population, which is why an
experiment that assumes the halo distribution will misread its own rate.}
\label{tab:velocities}
\begin{tabular}{@{}>{\RaggedRight\arraybackslash}p{0.44\textwidth}r@{}}
\toprule
\textbf{Population} & \textbf{Characteristic velocity} \\
\midrule
Galactic halo, typical & $\sim 220\ \kms$ \\
Local Sun's motion about the Galactic centre & $\sim 232\ \kms$ \\
Net wind speed at the Earth & $\sim 300\ \kms$ \\
Solar System escape velocity at 1~AU & $\sim 42\ \kms$ \\
Solar System escape velocity at the Sun & $\sim 72\ \kms$ \\
Captured population near the Earth & $\sim 10$--$40\ \kms$ \\
\bottomrule
\end{tabular}
\end{table}

\tabref{tab:velocities} is a table that a laboratory experimenter needs to
have in front of them, because the captured population and the halo
population differ by more than an order of magnitude in velocity, and the
recoil spectrum depends on the velocity through the kinematics of the
scattering. A limit computed for one population does not apply to the other.

\section{Why the enhancement is robust}The enhancement factor near the Sun depends
on the halo model only weakly, because it is set by the Solar System's own
gravitational field rather than by the halo's. The halo density sets the
incident flux; the capture cross section and the focusing geometry set the
captured distribution. This is why the enhancement is a robust prediction
while the absolute flux remains uncertain.

The consequence is that there is a region --- the inner Solar System --- where
the dark matter density is high enough to matter for a laboratory experiment,
and where the population's properties are determined by Solar System
dynamics rather than by the galaxy.

\section{The three measurements this enables}The bound population enables three
measurements that are not otherwise possible.

The \emph{local density measurement} is the most direct. If the bound
population can be distinguished from the unbound one --- by recoil energy, by
direction, or by the modulation with the Earth's position --- then the
enhancement factor becomes a measurement rather than a calculation, and the
halo model used in every laboratory limit can be tested directly.

The \emph{directional measurement} of the wind is the second. Its direction
reflects the Sun's motion in the galaxy and the Earth's position in its
orbit; its magnitude is a function of the halo velocity distribution, and
therefore of the quantity that enters every laboratory limit as the velocity
distribution. A directional measurement constrains the halo velocity
distribution at the Solar System, which propagates into the cross-section
limits.

The \emph{velocity distribution measurement} is the third, and it is the one
the dark matter hail proposal is designed for: because the scattered flux is
dominated by the slowest particles that can be deflected in, a measurement of
the flux measures the halo's speed distribution in a narrow band just above
the threshold, which is a sharper probe of that distribution than any bulk
measurement provides.

\section{Why this is under-used}Three reasons, and they are all practical.

The experiments are hard. A target thin enough to be sensitive to the
low-energy recoils of the bound population must be made of light elements,
must be radiopure, and must be shielded from the terrestrial cosmic-ray
background that is orders of magnitude larger than the signal.

The field does not have an instrument culture for this. The
gravitational-wave and neutrino communities have instruments built for
their own backgrounds; the dark matter community has instruments built for
underground backgrounds. A Solar System experiment requires neither
underground nor deep-space conditions, and it requires a different
engineering tradition from both.

And the connection to the exclusion curves is not obvious to people who
did not grow up in the field. The velocity distribution is usually treated
as a nuisance parameter; it is in fact a measurable quantity, and the
measurements that would constrain it are within reach.

\begin{keytakeaways}
\item The gravitationally captured population in the Solar System is
four-to-five orders of magnitude denser than the local halo and moves ten
times slower.
\item It enables three measurements nothing else can make: the local density,
the wind direction, and the halo velocity distribution.
\item The velocity distribution enters every laboratory limit as a nuisance
parameter, and the Solar System is the only place where it can be measured
directly.
\item The obstacle is engineering and instrumentation culture, not physics.
\end{keytakeaways}
